\begin{abstract}
Vehicle-to-Vehicle (V2V) communication is essential for safety-critical applications
such as collision avoidance and emergency braking in highway environments. However,
existing authentication schemes---including certificate-based, signature-based, and
recent zero-knowledge proof (ZKP) approaches---still face challenges in simultaneously
achieving ultra-low authentication latency, strong privacy guarantees, and scalability
under dense traffic conditions. Prior ZKP-based vehicular authentication schemes treat
each authentication as a full-cost independent operation and do not address
high-frequency proof generation, session-based credential management, or efficient
anonymous credential revocation suitable for high-speed V2V scenarios.

We propose ULP-V2V-Auth, an authentication framework combining three key design
elements: (1)~Anonymous Session Tokens (ASTs) with Merkle tree-based revocation,
decoupling identity from authentication and enabling $O(\log n)$ revocation overhead
versus $O(n)$ for CRL-based schemes; (2)~Groth16 zk-SNARK proofs with a
proof-slot cache model, where complete proofs are precomputed offline during AST
acquisition---each committed to a predicted BSM content---and only a Poseidon-2
message binding hash (0.44\,ms) is computed online at broadcast time, achieving
sub-1\,ms per-message authentication latency on standard vehicular On-Board Units
(OBUs); and (3)~a batch verification mechanism exploiting Groth16's algebraic
structure to reduce verification from $3k$ to $k+3$ pairing operations for $k$
simultaneous proofs, with a divide-and-conquer fallback (DCV) that bounds adversarial
batch-injection cost at $O(w \log k)$ sub-batch operations for $w$ invalid proofs.
Our implementation on BN254 (depth-16 circuit, $\sim$9{,}100 constraints) is directly
benchmarked on a Raspberry Pi~4 Model~B (ARM Cortex-A72, OBU-class proxy),
measuring: full prove 3{,}841\,ms with snarkjs and \textbf{1{,}753\,ms with rapidsnark}
($2.19\times$ speedup); batch verification 400\,ms for $k=30$ proofs
($\mathbf{4.65\times}$ measured speedup, exceeding the theoretical $2.73\times$
pairing-count ratio); and an online authentication cost of \textbf{0.439\,ms}
(0.44\% of the 100\,ms BSM cycle). We present a formal security analysis, a
quantitative comparison of Groth16 against PLONK and zk-STARK for V2V constraints,
a resilience analysis under real-world channel conditions including packet loss and
congestion, and a detailed experiment design with direct hardware benchmark results.
We further address open design challenges: secure AST issuance, event-driven AST
request policy, epoch-aware Merkle tree management, the Groth16 trusted setup
ceremony, concurrent sender/receiver operation, and proof cache sustainability.
\end{abstract}
